\section{反函数与可逆性}
\label{sec:02-Calculus/01-Functions-and-Precalculus/03}

\subsection{温度转换引发的追问}

摄氏度和华氏度的转换公式你见过很多次：

\[
F = \frac{9}{5}C + 32.
\]

朋友从纽约发消息说今天 68\(^\circ\)F。你把这数扔进公式反向求解：

\[
68 = \frac{9}{5}C + 32
\quad\Longrightarrow\quad
C = (68-32)\times\frac{5}{9} = 20.
\]

舒服。20\(^\circ\)C，刚好是舒服的室温。

但你做的不只是解方程——你\textbf{撤销}了一个函数。华氏度函数 \(F(C)\) 把摄氏度变成华氏度，而你找到了它的逆过程：把华氏度变回摄氏度。写成函数形式：

\[
F^{-1}(F) = \frac{5}{9}(F-32).
\]

验证一下：\(F^{-1}(68)=20\)，\(F^{-1}(F(20))=20\)。来回一趟，回到原点。

这个``撤销''动作在数学和现实中有无数分身：加密和解密，压缩和解压，微分和积分，前进和后退。但什么时候你能安全地撤销一个过程？什么时候撤销会出问题？

\subsection{碰撞：当信息在映射中丢失}

考虑另一个转换。某个物理实验里，能量 \(E\) 和温度 \(T\) 的关系是：

\[
E = T^2.
\]

如果探测器显示 \(E=100\)（焦耳），温度是多少？

\[
100 = T^2 \quad\Longrightarrow\quad T = 10 \;\text{或}\; T = -10.
\]

两个输入（\(10\) 和 \(-10\)）映射到同一个输出（\(100\)）。只看输出，你无法判断原始输入是哪一个——信息在平方的过程中丢失了。

这就是反函数存在的核心障碍：若函数把两个不同的输入送到同一个输出，仅看输出就不可能唯一恢复输入。代数上``解不出唯一解''只是这个结构问题的代数表现。

多试几组数感受这种丢失：

\[
\begin{aligned}
f(x)=x^2 &: \quad f(3)=9,\; f(-3)=9,\; f(0.5)=0.25,\; f(-0.5)=0.25.\\[4pt]
f(x)=\sin x &: \quad f(0)=0,\; f(\pi)=0,\; f(2\pi)=0,\; f(-\pi)=0.\\[4pt]
f(x)=|x| &: \quad f(5)=5,\; f(-5)=5.
\end{aligned}
\]

这些函数都有``碰撞''——不同输入撞到同一个输出。从输出往回找输入时，岔路口出现了。

\subsection{三种分类：单射、满射、双射}

为了精确描述一个函数是否可逆，需要三条判断标准。设 \(f:A\to B\)：

\textbf{单射（injective）}：不同的输入不会碰撞。正式地说：
\[
x_1\neq x_2 \;\Longrightarrow\; f(x_1)\neq f(x_2).
\]
等价地，若 \(f(x_1)=f(x_2)\) 则必 \(x_1=x_2\)。

在图像上，单射对应水平线检验：任意一条水平线 \(y=c\) 与图像至多相交一次。如果有一条水平线穿过图像两次，那这两个交点的横坐标就是一组碰撞的输入。（对照 \hyperref[sec:02-Calculus/01-Functions-and-Precalculus/01]{``从实数轴到函数''}一节的竖线检验：竖线检验判定``是否为函数''——一个输入不能有两个输出；水平线检验判定``是否可逆''——一个输出不能来自两个输入。两条检验分别守卫函数定义和可逆性的大门。）

\(f(x)=x^2\) 在 \(\R\) 上不是单射：水平线 \(y=4\) 穿过图像于 \(x=2\) 和 \(x=-2\) 两处。\(f(x)=x^3\) 在 \(\R\) 上是单射：任何水平线只穿过一次——立方函数严格上升，永不回头。

\textbf{满射（surjective）}：陪域 \(B\) 中的每一个元素都能被某个输入打到。即 \(f(A)=B\)——值域完全填满陪域，不留空隙。

\(f:\R\to\R,\; f(x)=x^2\) 不是满射：陪域包含 \(-1\)，但没有实数平方后等于 \(-1\)。但 \(f:\R\to[0,\infty),\; f(x)=x^2\) 是满射——陪域被限制到了实际能到达的范围。

\textbf{双射（bijective）}：同时满足单射和满射。这是``可逆''的精确条件——每个输出恰好对应唯一一个输入。

只有双射才存在反函数 \(f^{-1}:B\to A\)，满足：
\[
f^{-1}(f(x)) = x \quad (\forall x\in A),\qquad
f(f^{-1}(y)) = y \quad (\forall y\in B).
\]

两条等式各自的论域不同：第一条说``先从 \(A\) 出发经过 \(f\) 到 \(B\)，再用 \(f^{-1}\) 回来，回到原来的 \(x\)''。第二条说``从 \(B\) 出发经过 \(f^{-1}\) 到 \(A\)，再用 \(f\) 回去，回到原来的 \(y\)''。两者合在一起，完整表达了``执行再撤销''。

\subsection{限制定义域：主动放弃一些输入，换取可逆性}

\(f(x)=x^2\) 在 \(\R\) 上不可逆，但你并不是束手无策。如果你主动声明：``我只考虑非负温度''，那么 \(E=T^2\) 的定义域从 \(\R\) 缩小为 \([0,\infty)\)。

在这个受限的定义域上，不同非负数平方后确实不同——碰撞消失了。于是：
\[
f:[0,\infty)\to[0,\infty),\quad f(x)=x^2
\]
成为双射，反函数为 \(f^{-1}(y)=\sqrt{y}\)。

你也可以选择另一支：限制到 \((-\infty,0]\)，此时反函数为 \(f^{-1}(y)=-\sqrt{y}\)。同样可以撤销，但撤销的结果是非正数。

选择哪一支不是任意的——它取决于你的问题场景。温度在开尔文温标下天然非负，你选第一支。如果一个函数描述``向左侧的位移''，你选第二支。限制定义域不是在数学上偷懒，而是在声明``在我的问题里，只有这段输入有意义''。

图像上，反函数和原函数关于直线 \(y=x\) 对称。原因很直接：\(f\) 的图像由有序对 \((x, f(x))\) 组成，\(f^{-1}\) 的图像由 \((f(x), x)\) 组成——横纵坐标互换。几何上就是关于 \(y=x\) 做镜像反射。

拿 \(f(x)=x^2\) 在 \([0,\infty)\) 上的图像（右半支抛物线）和 \(f^{-1}(y)=\sqrt{y}\) 的图像放在一起，后者正是前者关于 \(y=x\) 的镜像——你看到的正是横轴和纵轴交换后的抛物线。

\subsection{记号陷阱：\(f^{-1}\) 不是 \(1/f\)}

写法的相似是一个历史遗留的坑：

\[
f^{-1}(x) \quad\text{表示}\quad \text{\(f\) 的反函数（撤销 \(f\)）},
\]
\[
f(x)^{-1} \;\text{或}\; \frac{1}{f(x)} \quad\text{表示}\quad \text{输出的乘法倒数}.
\]

用具体函数看清楚：

\[
f(x) = 2x+3.
\]

反函数：解 \(y=2x+3\) 得 \(x=\frac{y-3}{2}\)，所以
\[
f^{-1}(x) = \frac{x-3}{2}.
\]

倒数：
\[
\frac{1}{f(x)} = \frac{1}{2x+3}.
\]

完全不一样的东西。\(f^{-1}(5)=\frac{5-3}{2}=1\)，而 \(1/f(5)=1/13\approx 0.077\)。

三角函数的记号是重灾区：\(\sin^{-1}x\) 按国际惯例表示反正弦（arcsin），不表示 \(1/\sin x\)。\(1/\sin x\) 有专用记法 \(\csc x\)。但 \(\sin^2 x\) 又表示 \((\sin x)^2\)……这套记号的矛盾是历史演化的结果，没有优雅的解决方案。你只能靠上下文和复合验证来确认：如果看到 \(\sin^{-1}(\sin x)\)，它只能是反正弦——倒数解释会得出一个完全没有意义的量纲。

\subsection{量纲帮忙检查逆过程}

如果 \(f\) 把``秒''变成``米''（比如路程函数），那么 \(f^{-1}\) 必须把``米''变回``秒''。反函数的输入和输出类型正好交换。

这不是形式上的对称——它能帮你快速排除错误答案。假设你推导出的``反函数''把米变成了米/秒，你就知道推导过程出错了。

匀速运动是一个干净的演示：

\[
s(t) = s_0 + vt \qquad (v>0).
\]

位置 \(s\) 是时间 \(t\) 的函数。在 \(t=0\) 时位置为 \(s_0\)，之后以速度 \(v\) 前进。反函数：

\[
t(s) = \frac{s - s_0}{v}.
\]

输入是位置（米），输出是到达该位置的时间（秒）。量纲检查：分子是米，分母是米/秒，结果是秒——通过。

但如果运动折返了呢？设 \(s(t)=t^2-4t\)，在 \(t=2\) 处掉头。\(s=0\) 对应 \(t=0\) 和 \(t=4\) 两个时刻——碰撞。全程反函数不存在。

你可以在单调的一段上定义局部反函数：在 \([0,2]\) 上（前进阶段），反函数为 \(t_1(s)=2-\sqrt{4+s}\)；在 \([2,\infty)\) 上（折返阶段），反函数为 \(t_2(s)=2+\sqrt{4+s}\)。每一个局部的反函数都是良定义的，但它们各自只管辖一段运动。

\subsection{复合函数的逆：拆封要倒序}

设你要出门：先穿袜子（\(f\)），再穿鞋（\(g\)）。出门的完整动作为：
\[
h = g\circ f.
\]

回家脱的时候，你得先脱鞋（\(g^{-1}\)），再脱袜子（\(f^{-1}\)）：
\[
h^{-1} = f^{-1}\circ g^{-1}.
\]

用公式写：
\[
(g\circ f)^{-1} = f^{-1}\circ g^{-1}.
\]

顺序颠倒了。这不是一个需要死记的规则——想想穿鞋袜的场景就永远不会忘。

代数量纲也印证这个顺序。如果 \(f:\text{时间}\to\text{温度}\)，\(g:\text{温度}\to\text{电压}\)，那么 \(g\circ f:\text{时间}\to\text{电压}\)。逆回去：\((g\circ f)^{-1}:\text{电压}\to\text{时间}\)。先通过 \(g^{-1}\) 把电压变回温度（先撤销外层），再通过 \(f^{-1}\) 把温度变回时间（再撤销内层）。

这个逆序规则在后续的换元积分、坐标变换和矩阵求逆中会反复出现。现在把它内化成直觉就够了——撤销多步操作，永远从最后一步开始倒退。

\subsection{可逆性的代价与收获}

这一节的核心消息可以浓缩成一句话：函数可逆的充要条件是没有信息碰撞，且输出空间被完全覆盖。

你获得的工具：
\begin{itemize}
  \item 用单射/满射/双射精确判断一个函数的可逆性。
  \item 通过限制定义域，把不可逆的函数变出可逆的分支。
  \item 反函数的代数求解、图像对称和量纲检验。
  \item 复合函数的反函数需要逆序撤销。
\end{itemize}

但可逆性不是免费的。限制定义域意味着你主动放弃了定义域的另一半——你只能撤销``被保留的那段历史''。对于复杂的非单调函数，可能存在多个可逆分支，你需要根据问题场景选择正确的分支，而这个选择本身需要额外的信息（比如知道温度非负）。

另外，存在反函数不意味着反函数有闭式公式。\(f(x)=x^5+x+1\) 在 \(\R\) 上严格单调上升（导数恒正），因此是双射，反函数一定存在。但想把 \(y=x^5+x+1\) 解成 \(x=\text{某表达式}\)？没有初等的闭式解。存在性和可计算性是两回事。
